% from ega0-17-fr.tex
\oldpage[0\textsubscript{IV-1}]{51}
est de la forme $\mathfrak{m}\widehat{A}$, où $\mathfrak{m}$ est un idéal
maximal de $A$ contenant $\mathfrak{J}$, et que l'homomorphisme canonique
$A\to\widehat{A}$ donne un homomorphisme \emph{injectif}
$A_{\mathfrak{m}}\to\widehat{A}_{\mathfrak{n}}$, tel que la topologie
$\mathfrak{n}\widehat{A}_{\mathfrak{n}}$-préadique de
$\widehat{A}_{\mathfrak{n}}$ induise sur $A_{\mathfrak{m}}$ la topologie
$\mathfrak{m}A_{\mathfrak{m}}$-préadique, et que $A_{\mathfrak{m}}$ soit
dense dans $\widehat{A}_{\mathfrak{n}}$ pour la topologie
$\mathfrak{n}\widehat{A}_{\mathfrak{n}}$-préadique~; par suite les
\emph{complétés} des anneaux locaux noethériens $A_{\mathfrak{m}}$ et
$\widehat{A}_{\mathfrak{n}}$ sont canoniquement \emph{isomorphes}. Or, par
hypothèse $A_{\mathfrak{m}}$ est régulier, donc il en est de même de son
complété (17.1.5) et comme le complété de $\widehat{A}_{\mathfrak{n}}$ est
régulier, il en est de même de $\widehat{A}_{\mathfrak{n}}$ par (17.1.5).
\end{proof}

\begin{corollary}[17.3.9]
\label{0.17.3.9-fr}
Soit $A$ un anneau noethérien, quotient d'un anneau noethérien régulier $B$.
Si $C$ est une $A$-algèbre de type fini, tout anneau de fractions $S^{-1}C$
de $C$ est quotient d'un anneau régulier.
\end{corollary}

\begin{proof}
En effet, $C$ est quotient d'un anneau de polynômes $A[T_1,\ldots,T_n]$~;
comme $A=B/\mathfrak{J}$, où $\mathfrak{J}$ est un idéal de $B$, on a
$A[T_1,\ldots,T_n]=B[T_1,\ldots,T_n]/\mathfrak{J}B[T_1,\ldots,T_n]$~; en
vertu de (17.3.7), on peut donc se borner au cas où $C$ est quotient d'un
anneau régulier $B'$~; mais si $S'$ est l'image réciproque de $S$ dans $B'$,
$S'$ est une partie multiplicative de $B'$ et $C$ est un anneau quotient de
$S'^{-1}B'$, si bien que finalement tout revient à voir que lorsque $C$ est
\emph{régulier}, il en est de même de $S^{-1}C$, ce qu'on a vu en (17.3.6).
\end{proof}

L'importance des quotients d'anneaux réguliers tient entre autres à la
propriété précédente et au fait que ce sont des anneaux \emph{caténaires}
(16.5.12).

Tous les anneaux importants dans les applications à la Géométrie algébrique
sont des quotients d'anneaux réguliers.

% (continued in ega0-18-fr.tex, before its first page marker)
\section{Compléments sur les extensions d'algèbres}
\label{section:0.18-fr}

Ce paragraphe rassemble un certain nombre de constructions fonctorielles sur
les anneaux, qui seront utilisées de façon répétée aux \S\S~19 et 20~; il ne
contient aucun résultat non trivial.

\subsection{Images réciproques d'anneaux augmentés}
\label{subsection:0.18.1-fr}

\begin{env}[18.1.1]
\label{0.18.1.1-fr}
Étant donné un anneau $A$ (non nécessairement commutatif), la catégorie des
\emph{$A$-anneaux} a pour objets les couples $(B,\rho)$ formés d'un anneau
$B$ et d'un homomorphisme d'anneaux $\rho:A\to B$, pour morphismes (dits aussi
\emph{$A$-homomorphismes}) les homomorphismes d'anneaux $u:B\to C$ tels que,
si $\rho:A\to B$ et $\sigma:A\to C$ sont les homomorphismes (dits
\emph{structuraux}) définissant la structure de $A$-anneau sur $B$ et $C$
respectivement, le diagramme
\begin{equation}
\label{0.18.1.1.1-fr}
\tag{18.1.1.1}
  \xymatrix{
    B\ar[r]^u & C\\
    & A\ar[lu]^\rho\ar[u]_\sigma
  }
\end{equation}
soit commutatif. Le \emph{noyau} $\mathfrak{J}$ de $u$ est un idéal bilatère
qui est un $A$-\emph{bimodule} (pour $\rho$).

Si $f:A'\to A$ est un homomorphisme d'anneaux, pour tout $A$-anneau
$(B,\rho)$, $(B,\rho\circ f)$ est un $A'$-anneau, et si $u:B\to C$ est un
$A$-homomorphisme du $A$-anneau $(B,\rho)$
